\chapter{Metodología y Planificación del Proyecto}

\section{Análisis de metodologías de desarrollo}

La elección de una metodología de desarrollo condiciona la planificación, el seguimiento y la entrega del 
proyecto. Dado que este TFG se desarrolla de forma individual y con un alcance acotado a un MVP, la 
metodología debe priorizar la simplicidad operativa y la adaptación al cambio frente a la coordinación de 
equipos. A continuación se analizan las alternativas y se justifica la elegida:

\subsection{Metodologías tradicionales frente a ágiles}

Las \textbf{metodologías tradicionales} (modelo en cascada, modelo en V) estructuran el desarrolo en fases 
secuenciales —requisitos, diseño implementación, pruebas— en las que cada fase debe completarse antes de 
iniciar la siguiente. Exigen una especificación cerrada desde el inicio y penalizan los cambios tardíos, cuyo 
coste crece a lo largo del cilo.

Las \textbf{metodologías ágiles} sustituyen la secuencia rígida por un desarrollo iterativo e incremental: 
el producto se construye en ciclos cortos que entregan funcionalidad utilizable y permiten incorporar cambios 
de forma continua. Priorizan el software funcionando, la colaboracion y la respuesta al cambio sobre la 
documentación exhaustiva y la planificación cerrada.

Para este proyecto, el enfoque tradicional es inadecuado: el alcance no estaba completamente definido al inicio 
(las funcionalidades se han priorizado y reordenado durante el desarrollo) y un único desarrollador no obtiene 
beneficio del control formal entre las fases que justifica la cascada. El enfoque ágil se ajusta mejor a un 
desarrollo exploratorio y sujeto a revisión continua.

\subsection{Comparación de marcos ágiles}

Dentro del paradigma ágil, se consideran tres marcos relevantes:

\textbf{Scrum.} Organiza el trabajo en \textit{sprints} de duración fija, con roles definidos (Product Owner, 
Scrum Master, equipo de desarrollo) y ceremonias periódicas (planificación, \textit{daily}, revisión, 
retrospectiva). Está diseñado para equipos y su valor reside en la coordinación; sus roles y ceremonias 
resultan redundantes en un desarrollo individual.

\textbf{Kanban.} Gestiona el trabajo mediante un flujo continuo visualizado en un tablero de columnas (p. ej. 
\textit{Por hacer / En curso / Hecho}), limitando el trabajo en curso (WIP) para evitar la sobrecarga y los 
cuellos de botella. No impone roles, ceremonias ni iteraciones de duración fija: se adapta al ritmo de trabajo.

\textbf{Lean.} Más que un marco de gestión, es un conjunto de principios orientados a maximizar el valor eliminando 
el desperdicio (\textit{waste}): todo aquello que no aporta valor al producto final. Promueve construir solo lo 
necesario, validar pronto y posponer decisiones hasta el último momento responsable.

\begin{table}[htbp]
    \centering
    \caption{Comparación de marcos ágiles aplicables al proyecto.}
    \label{tab:comparacion-agiles}
    \begin{tabularx}{\textwidth}{@{}l *{3}{>{\RaggedRight\arraybackslash}X}@{}}
        \toprule
        \textbf{Criterio} & \textbf{Scrum} & \textbf{Kanban} & \textbf{Lean} \\
        \midrule
        Unidad de planificación      & Sprint (fijo)            & Flujo continuo & Flujo de valor \\
        Roles obligatorios           & Sí (3)                   & No             & No \\
        Ceremonias                   & Sí (4)                   & No             & No \\
        Idóneo para 1 desarrollador  & Bajo                     & Alto           & Alto \\
        Adaptación al cambio         & Media (intra-sprint rígido) & Alta        & Alta \\
        \bottomrule
    \end{tabularx}
\end{table}

\subsection{Justificación de la elección: Kanban Lean}

Se ha optado por un enfoque \textbf{Kanban con principios Lean} por las siguientes razones:

\textbf{Adecuación al desarrollo individual.} Scrum se descarta porque sus roles y ceremonias presuponen un equipo; 
mantenerlos en solitario añade sobrecargaa sin aportar coordinación. Kanban no impone esa estructura: el desarrollador
gestiona un único flujo de trabajo sin reuniones ni roles artificiales.

\textbf{Flujo continuo frente a iteraciones rígidas.} El alcance del MVP se ha reordenado durante el desarrollo 
(priorización de la establización de HU-1 a HU-4 frente a funcionalidades posteriores). El flujo continuo de Kanban absorbe
esta repriorización sin el coste de cerrar y replanificar \textit{sprints} de duración fija.

\textbf{Límite de WIP y foco.} Limitar el trabajo en curso obliga a completar una funcionalidad antes de iniciar la siguiente,
lo que en un desarrollo individual previene la dispersión y la acumulación de tareas a medio terminar.

\textbf{Principios Lean: minimizar el desperdicio.} La restricción de recursos propia de un TFG hace del principio Lean de 
eliminar el desperdicio una guía natural: se construye únicamente la funcionalidad que aporta valor al MVP, posponiendo 
deliberadamente todo lo accesorio (análisis emocional, persistencia, llamada de emergencia) hasta que el núcleo esté validado.
Esta decisión de acotar el alcance es una aplicación directa del pricnipio de \textit{construir solo lo necesario}.

La combinación resultante ---un tablero Kanban que visualiza el flujo, con WIP limitado, gobernado por el criterio de Lean
de priorizar valor y diferir lo accesorio--- ofrece el equilibrio entre disciplina y flexibilidad que requiere un desarrollo
individual, exploratorio y con plazo fijo.

\section{Aplicación de la metodología al proyecto}

Definidos los principios metodológicos en la sección anterior, este apartado documenta su aplicación concreta al desarrollo
de Quiet Companion: los actores implicados, la formalización de las Historias de Usuario como unidades de trabajo, el desglose 
técnico de las mismas, la estructura del tablero Kanban, el plan de iteraciones y el cronograma resultante.

\subsection{Definición de actores y roles}
Se distinguen dos planos de acores: los del \textbf{proyecto o producto} (quienes participan activamente en el desarrollo y 
la supervisión) y los del \textbf{sistema} (quienes interactúan con el producto desarrollado).

Actores del proyecto:
\begin{itemize}
    \item \textbf{Desarrollador Único.} Asume todos los roles técnicos: ingeniería de requisitos, diseño, implementación,
    pruebas, documentacion, y despliegue. Esta concentración de roles es lo que justifica la elección metodológica de 3.0.1
    \item \textbf{Tutores académicos.} Supervisan el progreso, validan decisiones de diseño y orientan el alcance. Las sesiones
    de tutoría actúan como puntos de control externos que sustituyten las ceremonias de Scrum (revisión y retrospectiva)
    que en un desarrollo individual carecen de sentido.
**Nota: también podría añadir el hecho de que he ido mostrando el avance del MVP a familiares y conocidos (personas mayores) como
"focus group" o muestra de target group, y también a otros compañeros programadores/informáticos para que me den feedback
\end{itemize}

Actores del sistema:
\begin{itemize}
    \item \textbf{Mayor independiente.} Persona mayor con autonomía cognitiva y técnica suficiente para gestionar la aplicación
    por sí misma. Accede al chat por voz y a los ajustes de accesibilidad en su propio dispositivo.
    \item \textbf{Mayor vinculado.} Persona mayor cuya configuración se delega en un familiar. Su interfaz queda restringida al 
    chat por voz, en el llamado \textit{Modo Ultra-Simple}. **Nota: Esto todavía habría que revisarlo
    \item \textbf{Familiar (companion).} Adulto que asume el rol de gestor de la configración del Mayor vinculado. Accede a un
    panel remoto desde su propio dispositvio y, en versiones futuras, a indicadores de bienestar emocional.
\end{itemize}

Estos tres actores del sistema se materializan en el modelo de datos mediante el campo \code{role} (\code{elder} | \code{companion}) y, en el caso 
del mayor, el campo \code{autonomy} (\code{independent} | \code{dependent}), lo que permite que una misma base de código sirva a los tres perfiles
 sin duplicación.

\subsection{Historias de Usuario como unidades de trabajo}

En la metodología adoptada, las Historias de Usuario cumplen un doble papel: son fuente de requisitos (se formalizan en el capítulo de especificación)
 y son unidades de trabajo del tablero Kanban. Este apartado las trata bajo el segundo enfoque.

El catálogo de Historias de Usuario gestionadas en el tablero es el siguiente:

\begin{table}[htbp]
    \centering
    \begin{tabularx}{\textwidth}{@{}c >{\RaggedRight\arraybackslash}X c c@{}}
        \toprule
        \textbf{ID} & \textbf{Título} & \textbf{Prioridad} & \textbf{Estado al cierre del MVP} \\
        \midrule
        HU-1 & Selección de Autonomía              & Alta  & Completada \\
        HU-2 & Interfaz Core Conversacional        & Alta  & Completada \\
        HU-3 & Interfaz Accesible y Comprensible   & Alta  & Completada \\
        HU-4 & Configuración Adaptativa            & Alta  & Completada \\
        HU-5 & Memoria Conversacional Básica       & Media & Diferida a trabajo futuro \\
        HU-6 & Seguimiento del Estado de Ánimo     & Media & Diferida a trabajo futuro \\
        \bottomrule
    \end{tabularx}
    \caption{Catálogo de Historias de Usuario gestionadas en el tablero Kanban.}
    \label{tab:catalogo-hu}
\end{table}

 Cada Historia de Usuario se ha gestionado como una tarjeta en el tablero, siguiendo una plantilla común que recoge la información 
mínima necesaria para entrar en flujo:

\begin{table}[htbp]
    \centering
    \begin{tabular}{|>{\raggedright\arraybackslash}p{0.48\textwidth}|>{\raggedright\arraybackslash}p{0.44\textwidth}|}
        \hline
        \textbf{Identificación}: HU-1 & Selección de Autonomía \\
        \hline
        \multicolumn{2}{|>{\raggedright\arraybackslash}p{0.94\textwidth}|}{%
            \textbf{Descripción}: \newline
            Como Mayor, quiero poder elegir, al abrir la aplicación por primera vez,
            si la gestionaré yo mismo o si me ayudará un familiar, para que el sistema
            me ofrezca la experiencia adecuada a mis capacidades tecnológicas.
        } \\
        \hline
        \textbf{Estimación}: L & \textbf{Prioridad}: 1 \\
        \hline
        \multicolumn{2}{|>{\raggedright\arraybackslash}p{0.94\textwidth}|}{%
            \textbf{Criterios de aceptación}:
            \begin{itemize}[nosep, leftmargin=*]
                \item La primera pantalla muestra dos opciones claras
                      («Usar por mi cuenta» / «Vincular con un familiar»).
                \item Si el usuario elige «Usar por mi cuenta», se habilita un botón
                      de «Ajustes» en su pantalla principal.
                \item Si elige «Vincular», se genera un código corto en pantalla y la
                      interfaz principal se bloquea hasta que el familiar se conecte.
            \end{itemize}
        } \\
        \hline
    \end{tabular}
    \caption{Historia de Usuario HU-1: Selección de Autonomía.}
    \label{tab:hu1}
\end{table}


**note: pasar la tabla a código latex

La aplicación efectiva del principio Lean de construir \textit{solo lo necesario} se manifiesta en la decisión deliberada 
de diferir HU-5 y HU-6: una vez constatado durante el desarrollo que la estabilización de HU-1 a HU-4 requería 
más esfuerzo del estimado inicialmente, se priorizó la entrega de un núcleo robusto sobre la cobertura amplia
 pero superficial de las seis historias. El detalle completo de cada HU, incluidos sus criterios de aceptación,
  se documenta en el capítulo de especificación.

\subsection{Tarjetas técnicas y desglose de tareas}
Cada Historia de Usuario se descompone en una o varias tarjetas técnicas, que son las unidades reales de trabajo
 del desarrollador. Esta separación responde al doble nivel de granularidad necesario: las HU expresan valor para
  el usuario; las tarjetas técnicas expresan trabajo concreto de ingeniería.

\begin{table}[htbp]
    \centering
    \begin{tabularx}{\textwidth}{@{}c c >{\RaggedRight\arraybackslash}X c@{}}
        \toprule
        \textbf{ID} & \textbf{HU} & \textbf{Título} & \textbf{Tipo} \\
        \midrule
        T-201 & HU-4 & Implementar \code{currentUserProvider} como \code{StreamProvider} único & Implementación \\
        T-202 & HU-4 & Migrar \code{TtsService} a inyección de dependencias vía Riverpod       & Refactor \\
        T-203 & HU-4 & Aplicar \code{textSizeMultiplier} en \code{home\_screen.dart}           & Implementación \\
        T-204 & HU-4 & Centralizar rutas de Firestore en \code{FirestorePaths}                & Refactor \\
        T-205 & HU-4 & Migrar \code{GeminiService} y \code{SpeechService} al patrón de DI      & Refactor \\
        T-206 & HU-4 & Romper acoplamiento horizontal \code{dashboard} $\rightarrow$ \code{onboarding} & Refactor \\
        \bottomrule
    \end{tabularx}
    \caption{Desglose de tarjetas técnicas de la HU-4 durante la fase de estabilización.}
    \label{tab:tarjetas-tecnicas}
\end{table}
\subsection{El tablero Kanban: columnas y flujo (TO DO / Doing / Done)}

\subsection{Plan de iteraciones / hitos (milestones)}

\begin{table}[htbp]
    \centering
    \begin{tabularx}{\textwidth}{@{}c >{\RaggedRight\arraybackslash}X c c@{}}
        \toprule
        \textbf{Hito} & \textbf{Objetivo} & \textbf{HU implicadas} & \textbf{Estado} \\
        \midrule
        M1 & Cimientos: Auth, modelo de datos, onboarding dual        & HU-1        & Completado \\
        M2 & Pipeline conversacional (STT $\rightarrow$ LLM $\rightarrow$ TTS) & HU-2, HU-3 & Completado \\
        M3 & Dashboard remoto y configuración adaptativa              & HU-4        & Completado \\
        M4 & Estabilización: refactors, DI, eliminación de deuda técnica & HU-1 a HU-4 & Completado \\
        M5 & Pulido de navegación y pruebas de usuario                & HU-1 a HU-4 & En curso \\
        M6 & Persistencia conversacional y memoria \textit{(futuro)}  & HU-5        & Diferido \\
        M7 & Análisis emocional y panel de bienestar \textit{(futuro)} & HU-6       & Diferido \\
        \bottomrule
    \end{tabularx}
    \caption{Plan de iteraciones e hitos del proyecto.}
    \label{tab:plan-hitos}
\end{table}

\subsection{Cronograma: diagrama de Gantt}
